% p.586
M step  M 步骤\quad 402, 405 \par
macrostate  宏观状态\quad 42 \par
MAE, see masked autoencoder  MAE, 参见“掩蔽自编码器” \par
MAF, see masked autoregressive flow  MAF, 参见“掩蔽自回归流” \par
Mahalanobis distance  马哈拉诺比斯距离\quad 60 \par
manifold  流形\quad 154, 445 \par
MAP, see maximum a posteriori  MAP, 参见“最大后验” \par
marginal probability  边缘概率\quad 25 \par
Markov blanket  马尔可夫毯\quad 297, 382 \par
Markov boundary, see Markov blanket  马尔可夫边界, 参见“马尔可夫毯” \par
Markov chain  马尔可夫链\quad 300, 377 \par
Markov chain Monte Carlo  马尔可夫链蒙特卡洛采样\quad 375 \par
Markov model  马尔可夫模型\quad 300 \par
Markov random field  马尔可夫随机场\quad 280 \par
masked attention  遮掩注意力\quad 328 \par
masked autoencoder  掩蔽自编码器\quad 483 \par
masked autoregressive flow  掩蔽自回归流\quad 472 \par
max-pooling  最大池化\quad 255 \par
max-unpooling  最大上采样\quad 272 \par
maximization step, see M step  最大化步骤, 参见“M步骤” \par
maximum a posteriori  最大后验\quad 49, 405 \par
maximum likelihood  最大似然\quad 34, 72, 101, 134 \par
MCMC, see Markov chain Monte Carlo  MCMC, 参见“马尔可夫链蒙特卡洛采样” \par
MDN, see mixture density network  MDN, 参见“混合密度网络” \par
mean  平均值\quad 32 \par
mean value theorem  中值定理\quad 43 \par
measure theory  测度论\quad 30 \par
mel spectrogram  梅尔频谱\quad 339 \par
message-passing  消息传递\quad 354 \par
message-passing neural network  消息传递神经网络\quad 354 \par
meta-learning  元学习\quad 165 \par
Metropolis algorithm  Metropolis 算法\quad 375 \par
Metropolis-Hastings algorithm  Metropolis-Hastings 算法\quad 378 \par
microstate  微观状态\quad 42 \par
mini-batches  小批量\quad 187 \par
minimum risk  最小风险\quad 127 \par
Minkowski loss  闵可夫斯基损失\quad 107 \par
missing at random  随机缺失\quad 406, 442 \par
missing data  缺失数据\quad 442 \par
mixing coefficient  混合系数\quad 75 \par
mixture component  混合情况\quad 75 \par
mixture density network  混合密度网络\quad 172 \par
mixture distribution  混合分布\quad 391 \par
mixture model  混合模型\quad 391 \par
mixture of Gaussians  高斯混合分布\quad 74, 173, 234, 397 \par
MLP, see multilayer perceptron  MLP, 参见“多层感知机” \par
MNIST data  MNIST 数据集\quad 421 \par
mode collapse  模式崩溃\quad 456 \par
model averaging  模型平均\quad 239 \par
model comparison  模型比较\quad 9 \par
model selection  模型选择\quad 12 \par
moment  矩\quad 33 \par
momentum  动量\quad 190 \par
Monte Carlo dropout  蒙特卡洛 dropout\quad 242 \par
Monte Carlo sampling  蒙特卡洛采样\quad 365 \par

% p.587
% 续上页：Monte Carlo sampling 的页码 365 载于本页首行
Moore-Penrose pseudo-inverse, see pseudo-inverse  摩尔-彭若斯伪逆, 参见“伪逆” \par
MRF, see Markov random field  MRF, 参见“马尔可夫随机场” \par
multi-class logistic regression  多类逻辑斯谛回归\quad 140 \par
multi-head attention  多头注意力\quad 313 \par
multilayer perceptron  多层感知机\quad 16, 149 \par
multimodal transformer  自注意力\quad 309 \par
multimodality  多模态\quad 173 \par
multinomial distribution  多项分布\quad 59, 82 \par
multiplicity  多重数\quad 42 \par
multitask learning  多任务学习\quad 165 \par
mutual information  互信息\quad 47 \par
n-gram model  n 元模型\quad 323 \par
naive Bayes model  朴素贝叶斯模型\quad 128, 294, 322 \par
nats  纳特\quad 41 \par
natural language processing  自然语言处理\quad 305 \par
natural parameter  自然参数\quad 81 \par
nearest-neighbours  最近邻\quad 88 \par
neocognitron  新认知机\quad 260 \par
Nesterov momentum  Nesterov 动量\quad 191 \par
neural ordinary differential equation  神经常微分方程\quad 472 \par
neuroscience  神经科学\quad 259 \par
NLP, see natural language processing  NLP, 参见“自然语言处理” \par
no free lunch theorem  无免费午餐定理\quad 221 \par
node  节点\quad 280, 349 \par
noise  噪声\quad 21 \par
noiseless coding theorem  无噪编码定理\quad 41 \par
noisy-OR  噪声 OR\quad 303 \par
non-identifiability  不可识别性\quad 437 \par
non-max suppression  非最大抑制\quad 269 \par
nonparametric methods  非参数化方法\quad 56, 85 \par
normal distribution, see Gaussian  正态分布, 参见“高斯分布” \par
normal equations  正规方程\quad 101 \par
normalized exponential, see softmax function  归一化指数, 参见“softmax 函数” \par
novelty detection  奇异值检测\quad 126 \par
object detection  目标检测\quad 265 \par
observed variable  观测变量\quad 287 \par
Old Faithful data  老忠实喷泉数据\quad 74 \par
on-hot encoding, see 1-of-K encoding  独热编码, 参见“1-of-K 编码” \par
one-shot learning  单样本学习\quad 165 \par
one-versus-one classifier  一对一分类器\quad 118 \par
one-versus-the-rest classifier  一对多分类器\quad 117 \par
online gradient descent  在线梯度下降\quad 186 \par
online learning  在线学习\quad 102 \par
ordered over-relaxation  有序过度放松\quad 382 \par
outer product approximation  外积近似法\quad 210 \par
outlier  离群值\quad 120, 126, 143 \par
over-fitting  过拟合\quad 9, 108, 400 \par
over-relaxation  过度放松\quad 382 \par
over-smoothing  过度平滑\quad 362 \par
padding  填充\quad 252 \par
parameter sharing  参数共享\quad 234, 284 \par

% p.588
parameter shrinkage  参数收缩\quad 103 \par
parameter tying, see parameter sharing  参数捆绑, 参见“参数共享” \par
parent node  父节点\quad 214, 281 \par
partition function  配分函数\quad 384 \par
Parzen estimator, see kernel density estimator  Parzen 估计器, 参见“核密度估计器” \par
Parzen window  Parzen 窗\quad 87 \par
PCA, see principal component analysis  PCA, 参见“主成分分析” \par
perceptron  感知机\quad 15 \par
periodic variables  周期变量\quad 76 \par
permutation matrix  排列矩阵\quad 350 \par
PixelCNN  PixelCNN\quad 338 \par
PixelRNN  PixelRNN\quad 338 \par
plate  板块\quad 286 \par
polynomial curve fitting  多项式曲线拟合\quad 55 \par
pooling  池化\quad 255 \par
positional encoding  位置编码\quad 317 \par
positive definite covariance  正定协方差矩阵\quad 61 \par
positive definite matrix  正定矩阵\quad 523 \par
posterior collapse  后验崩塌\quad 491 \par
posterior probability  后验概率\quad 28 \par
power method  幂方法\quad 424 \par
pre-activation  预激活\quad 15 \par
pre-processing  预处理\quad 17 \par
pre-training  预训练\quad 164, 334 \par
precision matrix  精度矩阵\quad 65 \par
precision parameter  精度\quad 32 \par
predictive distribution  预测分布\quad 37, 105 \par
prefix prompt  前缀提示\quad 335 \par
principal component analysis  主成分分析\quad 422, 430, 481 \par
principal subspace  主子空间\quad 422 \par
prior  先验\quad 227 \par
prior knowledge  先验知识\quad 17, 220 \par
prior probability  先验概率\quad 28, 126 \par
probabilistic graphical model, see graphical model  概率图模型, 参见“图模型” \par
probabilistic PCA  概率 PCA\quad 430 \par
probability  概率\quad 23 \par
probability density  概率密度\quad 28 \par
probability theory  概率论\quad 22 \par
probit function  probit 函数\quad 142 \par
probit regression  probit 回归\quad 141 \par
product rule of probability  概率的乘积法则\quad 24, 25, 280 \par
prompt  提示\quad 335, 511 \par
prompt engineering  提示工程\quad 335 \par
proposal distribution  提议分布\quad 369, 371, 375 \par
pseudo-inverse  伪逆\quad 101, 120 \par
pseudo-random numbers  伪随机数\quad 366 \par
quadratic discriminant  二次判别式\quad 133 \par
radial basis functions  径向基函数\quad 155 \par
random variable  随机变量\quad 24 \par
raster scan  栅格扫描\quad 338 \par
readout layer  读出层\quad 357 \par
real NVP normalizing flow  实数 NVP 标准化流\quad 468 \par
receiver operating characteristic, see ROC curve  受试者工作特征曲线, 参见“ROC 曲线” \par
receptive field  感受野\quad 250, 355 \par
recurrent neural network  递归神经网络\quad 324 \par
regression  回归\quad 3 \par
regression function  回归函数\quad 106 \par
regularization  正则化\quad 11, 219 \par
regularized least squares  正则化最小二乘法\quad 103 \par

% p.589
reject option  拒绝选项\quad 125, 127 \par
rejection sampling  拒绝采样\quad 369 \par
relative entropy  相对熵\quad 45 \par
reparameterization trick  重参数化技巧\quad 488 \par
representation learning  表示学习\quad 19, 163 \par
residual block  残差块\quad 237 \par
residual connection  残差连接\quad 19, 236 \par
residual network  残差网络\quad 237 \par
resnet, see residual network  ResNet, 参见“残差网络” \par
responsibility  责任\quad 75, 398 \par
RLHF  RLHF\quad 335 \par
RMS error, see root-mean-square error  RMS 误差, 参见“均方根误差” \par
RMSProp  RMSProp\quad 193 \par
RNN, see recurrent neural network  RNN, 参见“递归神经网络” \par
robot arm  机械臂\quad 172 \par
robustness  鲁棒性\quad 120 \par
ROC curve  ROC 曲线\quad 129 \par
root-mean-square error  均方根误差\quad 9 \par
saliency map  显著性图\quad 262 \par
same convolution  等大卷积\quad 253 \par
sample mean  样本均值\quad 34 \par
sample variance  样本方差\quad 34 \par
sampling  采样\quad 365 \par
sampling-importance-resampling  采样-重要性-重采样\quad 373 \par
scale invariance  尺度不变性\quad 222 \par
scaled self-attention  缩放自注意力\quad 312 \par
scaling hypothesis  缩放假设\quad 306 \par
Schur complement  舒尔补\quad 67 \par
score function  得分函数\quad 386, 505 \par
score matching  得分匹配\quad 505 \par
self-attention  自注意力\quad 309 \par
self-supervised learning  自监督学习\quad 5, 320 \par
semi-supervised learning  半监督学习\quad 358 \par
sequential estimation  序贯估计\quad 73 \par
sequential gradient descent  序贯梯度下降\quad 103 \par
sequential learning  序贯学习\quad 102 \par
SGD, see stochastic gradient descent  SGD, 参见“随机梯度下降” \par
shared parameters, see parameter sharing  共享参数, 参见“参数共享” \par
shared weights  共享权重\quad 251 \par
shattered gradients  破碎梯度\quad 236 \par
shrinkage  收缩\quad 11 \par
sigmoid, see logistic sigmoid  sigmoid, 参见“逻辑斯谛 sigmoid” \par
singular value decomposition  奇异值分解\quad 102 \par
SIR, see sampling-importance-resampling  SIR, 参见“采样-重要性-重采样” \par
skip-grams  跳字\quad 320 \par
skip-layer connections  跳层连接\quad 237 \par
sliding window  滑动窗口\quad 267 \par
smoothing parameter  平滑参数\quad 95 \par
soft ReLU  soft ReLU\quad 161 \par
soft weight sharing  软权重共享\quad 234 \par
softmax function  softmax 函数\quad 83, 132, 171, 174, 310 \par
softplus activation function  softplus 激活函数\quad 160 \par
sparse autoencoders  稀疏自编码器\quad 482 \par
sparse connections  稀疏连接\quad 251 \par
sparsity  稀疏\quad 229 \par
sphering  球化\quad 429 \par
standard deviation  标准差\quad 32 \par
standardizing  标准化\quad 394, 428 \par
state-space model  状态空间模型\quad 301 \par
statistical bias, see bias  统计偏差, 参见“偏差” \par

% p.590
% 续上页：statistical bias 的“偏差”一条续于本页首行
statistical independence, see independent variables  统计独立, 参见“独立变量” \par
steepest descent  最速下降\quad 185 \par
Stein score, see score function  斯坦因得分, 参见“得分函数” \par
Stirling's approximation  斯特林公式\quad 42 \par
stochastic  随机\quad 7 \par
stochastic differential equation  随机微分方程\quad 508 \par
stochastic gradient descent  随机梯度下降\quad 17, 186 \par
stochastic variable  随机变量\quad 26 \par
strided convolution  跨步卷积\quad 253 \par
strides  跨步\quad 267 \par
structured data  结构化数据\quad 247, 365 \par
style transfer  风格迁移\quad 274 \par
sufficient statistics  充分统计量\quad 57, 58, 72, 84 \par
sum rule of probability  概率的加和法则\quad 24, 25, 280 \par
sum-of-squares error  平方和误差函数\quad 8, 37, 119 \par
supervised learning  监督学习\quad 3, 358 \par
support vector machine  支持向量机\quad 156 \par
SVD, see singular value decomposition  SVD, 参见“奇异值分解” \par
SVM, see support vector machine  SVM, 参见“支持向量机” \par
swish activation function  swish 激活函数\quad 179 \par
symmetry  对称性\quad 222 \par
symmetry breaking  对称性破坏\quad 188 \par
tail-to-tail path  尾尾相接\quad 290 \par
tangent propagation  切线传播\quad 223 \par
temperature  温度\quad 330 \par
tensor  张量\quad 168, 253 \par
test set  测试集\quad 9, 13 \par
text-to-speech  文本语音转换\quad 340 \par
tied parameters, see parameter sharing  捆绑参数, 参见“参数共享” \par
token  token\quad 308 \par
tokenization  分词\quad 321 \par
training set  训练集\quad 3 \par
transductive  直推式\quad 349, 358 \par
transductive learning  直推式学习\quad 358 \par
transfer learning  迁移学习\quad 3, 164, 188, 331 \par
transformers  Transformer\quad 305 \par
transition probability  转移概率\quad 377 \par
translation invariance  平移不变性\quad 222 \par
transpose convolution  转置卷积\quad 272 \par
tri-gram model  三元模型\quad 323 \par
TTS, see text-to-speech  TTS, 参见“文本语音转换” \par
U-net  U-Net\quad 273 \par
undetermined multiplier, see Lagrange multiplier  未定乘子, 参见“拉格朗日乘子” \par
undirected graphical model  无向图模型\quad 280 \par
uniquenesses  唯一性\quad 436 \par
universal approximation theorems  通用近似器\quad 159 \par
unobserved variable, see latent variable  未观测变量, 参见“隐变量” \par
unsupervised learning  无监督学习\quad 4, 163 \par
utility function  效用函数\quad 124 \par
VAE, see variational autoencoder  AE, 参见“变分自编码器” \par
valid convolution  有效卷积\quad 253 \par
validation set  验证集\quad 13 \par

% p.591
vanishing gradient  梯度消失\quad 196, 325 \par
variance  方差\quad 31, 32, 109 \par
variational autoencoder  变分自编码器\quad 484 \par
variational inference  变分推断\quad 412 \par
variational lower bound, see evidence lower bound  变分下界, 参见“证据下界” \par
vector quantization  向量量化\quad 338, 396 \par
vertex, see node  顶点, 参见“点” \par
vision transformer  视觉 Transformer\quad 336 \par
von Mises distribution  冯·米塞斯分布\quad 76 \par
voxel  体素\quad 249 \par
Wasserstein distance  Wasserstein 距离\quad 457 \par
Wasserstein GAN  Wasserstein GAN\quad 458 \par
wavelets  小波\quad 99 \par
weakly supervised  弱监督\quad 167 \par
weight decay  权重衰减\quad 11, 225 \par
weight parameter  权重参数\quad 15, 156 \par
weight sharing, see parameter sharing  权重共享, 参见“参数共享” \par
weight vector  权重向量\quad 116 \par
weight-space symmetry  权重空间的对称性\quad 161 \par
WGAN, see Wasserstein GAN  WGAN, 参见“Wasserstein GAN” \par
whitening  白化\quad 429 \par
Woodbury identity  伍德伯里恒等式\quad 518 \par
word embedding  词嵌入\quad 320 \par
word2vec  word2vec\quad 320 \par
wrapped distribution  环绕分布\quad 80 \par
Yellowstone National Park  黄石国家公园\quad 74 \par

% p.592 blank
